\subsection{Aufbau der Arbeit}
\label{subsec:01-05-02_structure}

Der in Abschnitt~\ref{subsec:01-05-01_sci-approach} beschriebene wissenschaftliche Ansatz spiegelt sich im Aufbau der Arbeit und in der Kapitelstruktur wider. Eine schematische Übersicht über die Kapitel und deren Beitrag zur Beantwortung der Forschungsfragen bietet Abbildung~\ref{TODO:fig:aufbau-der-arbeit}, in der der inhaltliche Verlauf von den Grundlagen über die empirische Analyse hin zu Diskussion und Handlungsempfehlungen visualisiert wird.

\hyperref[ch:01_purpose-and-objectives]{Kapitel~1 \enquote{Zweck und Ziele der Arbeit}} führt in die Problemstellung und den Zweck der Untersuchung ein, beschreibt den Projektkontext und die Relevanz, leitet Zielsetzung und Forschungsfragen ab, grenzt den Untersuchungsrahmen ein und erläutert den wissenschaftlichen Ansatz sowie den Aufbau der Arbeit. Damit schafft das Kapitel die inhaltliche und methodische Grundlage für die folgenden Ausführungen.

\hyperref[ch:02_theoretical-foundations]{Kapitel~2 \enquote{Grundlagen}} stellt die theoretische Basis der Arbeit bereit. Es definiert zentrale Begriffe, führt relevante UX-Methoden für interne Plattformen ein, beschreibt Informationsarchitektur als kognitive Orientierungshilfe und beleuchtet den Zusammenhang von UX und digitaler Transformation. Die dort entwickelten Konzepte und Modelle bilden den theoretischen Bezugsrahmen zur Beantwortung der Forschungsfragen RQ-1 bis RQ-3.

\hyperref[ch:03_pre-analysis]{Kapitel~3 \enquote{Voranalyse}} widmet sich der detaillierten Betrachtung des Untersuchungsfelds. Es beschreibt den Kontext der betrachteten Organisation, charakterisiert die bestehende Plattform, analysiert den Ausgangszustand von Informationsarchitektur und UX und identifiziert maßgebliche Nutzergruppen und Stakeholder. Zudem werden vorhandene UX-Strategien, zentrale Herausforderungen und Potenziale herausgearbeitet und daraus der Forschungsfokus inklusive Annahmen abgeleitet, was vor allem zur empirischen Fundierung von RQ-1 und RQ-2 beiträgt.

\hyperref[ch:04_methodology]{Kapitel~4 \enquote{Forschungsdesign und Methodik}} konkretisiert den wissenschaftlichen Ansatz. Es beschreibt den Forschungsansatz der qualitativen Einzelfallstudie, die Datenerhebung mittels Experteninterview und Usability-Tests, die Auswahl der Interviewpartner und Testpersonen, die Leitfadenentwicklung sowie Durchführung und Dokumentation. Darüber hinaus werden Vorgehen und Kriterien der qualitativen Inhaltsanalyse sowie Gütekriterien, Limitationen und Abgrenzungen diskutiert. Dieses Kapitel zeigt damit transparent, wie die Daten zur Beantwortung der Forschungsfragen erhoben und ausgewertet werden.

\hyperref[ch:05_analysis-and-results]{Kapitel~5 \enquote{Analyse und Ergebnisse}} präsentiert die empirischen Resultate der Arbeit. Es gibt zunächst einen Überblick über die erhobenen Daten, analysiert Wahrnehmungen und Herausforderungen in der bestehenden Informationsarchitektur und untersucht die Rolle von Personas und UX-Strategien in der Praxis. Anschließend werden identifizierte Erfolgsfaktoren und Muster sowie Akzeptanzfaktoren im organisationalen Kontext herausgearbeitet. Die Ergebnisse dieses Kapitels sind zentral für die spätere Beantwortung aller drei Forschungsfragen.

\hyperref[ch:06_discussion]{Kapitel~6 \enquote{Diskussion}} ordnet die gewonnenen Ergebnisse in den theoretischen Kontext ein. Es verknüpft die empirischen Befunde mit den in \hyperref[ch:02_theoretical-foundations]{Kapitel~2} dargestellten Konzepten, beantwortet die Forschungsfragen RQ-1 bis RQ-3 explizit und leitet Implikationen für UX-Strategien in Organisationen ab. Zudem werden Potenziale für zukünftige Plattformentwicklungen diskutiert, wodurch das Kapitel die Brücke zwischen Fallstudie und allgemeiner wissenschaftlicher Einordnung schlägt.

\hyperref[ch:07_recommendations]{Kapitel~7 \enquote{Handlungsempfehlungen für die Praxis}} bündelt die praxisorientierten Konsequenzen der Arbeit. Es formuliert Strategien zur nutzerzentrierten Informationsarchitektur, beschreibt Ansätze zur Integration von UX-Methoden in Entwicklungs- und Transformationsprozesse und skizziert Kommunikations- und Change-Maßnahmen zur Förderung der Akzeptanz interner Plattformen. Die Empfehlungen stützen sich direkt auf die Ergebnisse und Diskussionen der vorangegangenen Kapitel.

Abschließend fasst \hyperref[ch:08_conclusion]{Kapitel~8 \enquote{Fazit und Ausblick}} die zentralen Erkenntnisse der Arbeit zusammen, reflektiert deren Grenzen und diskutiert verbleibende Forschungs- und Entwicklungsperspektiven. Damit schließt das Kapitel den Bogen zurück zu den in \hyperref[ch:01_purpose-and-objectives]{Kapitel~1} formulierten Zielen und Forschungsfragen und zeigt auf, in welchen Bereichen weiterführende Untersuchungen sinnvoll erscheinen.